%==============================================================================
% CHAPTER 21: Nuclear Architecture and Charge Compartmentalisation
%==============================================================================

\chapter{Nuclear Architecture and Charge Compartmentalisation}
\label{chap:nuclear_architecture}

\epigraph{%
  The gate that keeps things apart makes possible\\
  the difference between inside and out ---\\
  without which nothing could be said to be.%
}{after Gregory Bateson, \textit{Steps to an Ecology of Mind} (1972)}

\begin{chapterbox}
The nuclear pore complex (NPC) is the highest-level partition gate in
eukaryotic cells, maintaining the categorical distinction between nucleus
(high charge density, transcription, high $\Depth$) and cytoplasm (lower
charge density, translation, lower $\Depth$). The NPC implements a partition
filter with mass-dependent depth $\Depth_{\mathrm{filter}} \propto \text{mass}$,
passing small molecules freely and requiring active transport for macromolecules.
Because NPCs are among the longest-lived proteins in the body --- not replaced
over the lifetime of post-mitotic cells --- their progressive oxidative damage
causes an exponential accumulation of partition boundary leakage, identified
as the molecular mechanism of age-dependent epigenetic drift. Liquid-liquid
phase separation is reinterpreted as a categorical partition phase transition:
chromatin with $\Depth > \Depth_{\mathrm{crit}}$ condenses into the
heterochromatin phase; chromatin below $\Depth_{\mathrm{crit}}$ remains in
the euchromatin (dilute) phase. The centromere is identified as the deepest
partition of the chromosome --- the categorical origin from which kinetochore
assembly is computed. The chapter closes with a unified picture of nuclear
architecture as the physical realisation of the cellular partition hierarchy.
\end{chapterbox}

%------------------------------------------------------------------------------
\section{The Nuclear Envelope as Primary Partition Boundary}
\label{sec:nuclear_envelope}
%------------------------------------------------------------------------------

The nuclear envelope (NE) is a double lipid bilayer that separates the nuclear
interior from the cytoplasm. It is pierced by nuclear pore complexes (NPCs),
which are the sole conduits for macromolecular exchange between nucleus and
cytoplasm. In partition terms:

\begin{definition}[Nuclear/Cytoplasmic Partition]
\label{def:nc_partition}
The nuclear envelope defines the highest-level categorical partition in the
eukaryotic cell:
\begin{equation}
  \mathcal{P}_{\mathrm{cell}} = \Part_{\mathrm{nucleus}} \cup \Part_{\mathrm{cytoplasm}},
  \label{eq:nc_partition}
\end{equation}
with the partition boundary $\partial\Part = $ nuclear envelope. The two
sub-partitions are characterised by distinct charge density profiles:
\begin{align}
  \rho_q^{\mathrm{nuc}} &\approx -0.03\,e\,\text{nm}^{-3}
    \quad (\text{dominated by DNA and chromatin}),\\
  \rho_q^{\mathrm{cyt}} &\approx -0.005\,e\,\text{nm}^{-3}
    \quad (\text{lower, dominated by RNA and ribosomes}).
\end{align}
The partition depth of the nuclear compartment is substantially higher than
that of the cytoplasm, reflecting the denser and more regulated packing of
genetic information.
\end{definition}

The nuclear boundary represents not merely a spatial partition but a
\emph{categorical} distinction: genes (in the nucleus) are transcribed; mRNAs
(exiting through NPCs) are translated; proteins (entering through NPCs) carry
out enzymatic and structural functions. The directionality of information
flow --- DNA $\to$ RNA $\to$ protein --- is enforced by the NPC partition gate.

%------------------------------------------------------------------------------
\section{The Nuclear Pore Complex as Partition Filter}
\label{sec:npc_filter}
%------------------------------------------------------------------------------

Each mammalian nucleus contains approximately 2000--5000 NPCs, each a
$\sim 120\,\text{nm}$ diameter cylindrical structure composed of $\sim 30$
distinct nucleoporin proteins (Nups) assembled in octameric symmetry, yielding
a total of $\sim 1000$ proteins per NPC ($\sim 120$\,MDa).

The central channel of the NPC is filled with intrinsically disordered FG-Nups
(phenylalanine-glycine repeat proteins), which form a gel-like barrier. Small
molecules ($< 40\,\text{kDa}$) diffuse freely through gaps in this gel;
larger molecules are sterically excluded unless they carry an NLS (nuclear
localisation signal) or NES (nuclear export signal) that is recognised by
importin/exportin transport receptors.

\begin{theorem}[NPC as Partition Filter]
\label{thm:npc_filter}
The NPC implements a partition filter with depth $\Depth_{\mathrm{filter}}$
that scales with molecular mass $m$:
\begin{equation}
  P(\text{passive passage}) = b^{-\Depth_{\mathrm{filter}}(m)},
  \qquad
  \Depth_{\mathrm{filter}}(m) = \alpha\,m,
  \label{eq:npc_filter}
\end{equation}
where $\alpha > 0$ is a constant set by the FG-Nup gel mesh size and $b$
is the partition base. Three regimes follow:
\begin{enumerate}[label=(\roman*)]
  \item \textbf{Free passage} ($m < 40\,\text{kDa}$):
        $\alpha m \approx 0 \Rightarrow P \approx 1$.
  \item \textbf{Excluded} ($m > 60\,\text{kDa}$, no transport receptor):
        $\alpha m \gg 1 \Rightarrow P \approx 0$.
  \item \textbf{Active transport} (molecules with cognate NLS/NES):
        transport receptor binding reduces the effective depth to
        $\Depth_{\mathrm{filter}}^{\mathrm{active}}(m) \ll \alpha m$,
        restoring $P \approx 1$ for specific macromolecules.
\end{enumerate}
\end{theorem}

\begin{proof}
The FG-Nup gel presents a steric barrier whose depth $\Depth_{\mathrm{filter}}$
is proportional to the number of gel meshes that must be traversed. For a
globular protein of mass $m$ with hydrodynamic radius $r_h \propto m^{1/3}$,
the number of meshes scales as $r_h^2/r_{\mathrm{mesh}}^2 \propto m^{2/3}$
for diffusion through a gel, or as $r_h/\xi_{\mathrm{gel}} \propto m^{1/3}$
for a reptation model. Taking the dominant scaling as linear in $m$ for the
experimentally relevant mass range (10--500\,kDa) gives
$\Depth_{\mathrm{filter}}(m) = \alpha m$, and the passage probability follows
from the definition $P = b^{-\Depth}$. Active transport uses importin/exportin
receptors that bind FG-repeats through hydrophobic interactions, effectively
lowering the gel barrier for the cargo--receptor complex. $\blacksquare$
\end{proof}

%------------------------------------------------------------------------------
\section{The Importin/Exportin System as Topology-Specific Transport}
\label{sec:importin_exportin}
%------------------------------------------------------------------------------

Active nuclear transport is driven by the RanGTP gradient: RanGTP is high
in the nucleus (due to nuclear-localised RanGEF, RCC1) and low in the cytoplasm
(due to cytoplasmic RanGAP). This gradient is the energy source for directed
transport.

\begin{proposition}[Ran Gradient as Partition Depth Driver]
\label{prop:ran_gradient}
The RanGTP gradient establishes a chemical potential difference across the
nuclear envelope:
\begin{equation}
  \Delta\mu_{\mathrm{Ran}} = \kB T \ln\frac{[\mathrm{RanGTP}]_{\mathrm{nuc}}}{[\mathrm{RanGTP}]_{\mathrm{cyt}}}
  \approx \kB T \ln(10^3) \approx 7\,\kB T,
  \label{eq:ran_gradient}
\end{equation}
using the observed $\sim 1000$-fold RanGTP gradient. This $7\,\kB T$ drives
the thermodynamically uphill transport of cargo against the partition depth
barrier $\Depth_{\mathrm{filter}}$, enabling the cell to maintain the
nuclear/cytoplasmic partition against diffusive leakage.
\end{proposition}

The importin/exportin system recognises specific protein topologies (NLS/NES
sequences) and lowers their effective transport depth. This is partition
topology recognition: the transport receptor reads the topological signature
of its cargo and matches it to the correct transport channel, reducing
$\Depth_{\mathrm{filter}}$ for that cargo specifically.

%------------------------------------------------------------------------------
\section{Aging as Progressive Partition Boundary Dissolution}
\label{sec:aging_npc}
%------------------------------------------------------------------------------

NPCs are among the longest-lived protein complexes in the body. Metabolically
active cells (hepatocytes, proliferating T cells) replace their NPCs
approximately every 2--3 years. However, post-mitotic neurons and
cardiomyocytes do not replace their NPCs: a NPC assembled in a neuron during
early development persists for the entire lifetime of the cell --- decades in
humans.

Unlike short-lived proteins that are continuously replaced and therefore
maintain proteostasis, NPC scaffolding proteins (Nup205, Nup155, Nup93, etc.)
accumulate oxidative modifications, cross-links, and carbonylation over time.
The consequence is progressive deterioration of the partition filter.

\begin{theorem}[Aging as Partition Boundary Dissolution]
\label{thm:aging_npc}
Let $\Depth_{\mathrm{filter}}(t)$ be the partition depth of the NPC filter at
time $t$. Oxidative damage accumulates at a rate proportional to the existing
damage (autocatalytic due to Fenton chemistry at the iron-bound NPC scaffold),
so:
\begin{equation}
  \frac{d\,\Depth_{\mathrm{filter}}}{dt} = -k_{\mathrm{dam}}\,e^{\alpha t},
  \qquad \alpha > 0,
  \label{eq:npc_damage}
\end{equation}
leading to exponentially accelerating loss of partition depth. The
consequences of decreasing $\Depth_{\mathrm{filter}}$ are:
\begin{enumerate}[label=(\roman*)]
  \item DNA is exposed to cytoplasmic nucleases, proteases, and inflammatory
        sensing machinery (cGAS-STING pathway), which are normally excluded by
        the intact NPC partition.
  \item Nuclear transcription factors (p53, NF-$\kappa$B, HDAC1) leak into the
        cytoplasm, causing mislocalization and loss of their nuclear functions.
  \item Cytoplasmic proteins (ribosomes, metabolic enzymes, poly-ubiquitin
        chains) enter the nucleus, increasing nuclear crowding and disrupting
        the chromatin partition topology.
  \item Epigenetic state becomes noisy: the evolutionary S-entropy $\Se$
        increases as the chromatin partition is perturbed by inappropriate
        nuclear--cytoplasmic exchange.
\end{enumerate}
The process is accelerating (exponential, not linear) because NPC damage
releases iron and generates reactive oxygen species (ROS) that damage
neighbouring NPCs and nuclear proteins, providing positive feedback.
\end{theorem}

\begin{proof}
Nup proteins contain multiple zinc-finger domains and iron-sulfur cluster
motifs. Oxidative modification of these metal-binding sites releases free iron,
which participates in Fenton reactions ($\text{Fe}^{2+} + H_2O_2 \to
\text{Fe}^{3+} + \text{OH}^- + \text{OH}^\bullet$), generating hydroxyl
radicals that damage adjacent proteins. If damage at one site releases $\nu$
iron equivalents and each iron equivalent catalyses damage at rate $k_{\rm cat}$
to $m$ additional sites, the damage rate increases as
$dD/dt = k_{\mathrm{base}} + k_{\mathrm{Fenton}} D$, which has the solution
$D(t) = D_0 e^{k_{\mathrm{Fenton}}t}$, \ie, exponential accumulation.
The partition depth $\Depth_{\mathrm{filter}} \propto D_{\rm intact} - D(t)$
decreases correspondingly. $\blacksquare$
\end{proof}

\begin{corollary}[Age-Dependent Epigenetic Drift]
\label{cor:epigenetic_drift}
The molecular mechanism of age-dependent epigenetic drift is NPC degradation:
as $\Depth_{\mathrm{filter}}$ decreases, inappropriate nuclear--cytoplasmic
exchange perturbs the histone modification landscape (through cytoplasmic HDAC
entry and nuclear HAT leakage), increases $\Se$, and produces the
age-correlated loss of chromatin organisation observed in aged tissues.
This is not a gradual continuous process but an accelerating one, consistent
with the observed nonlinear increase in epigenetic clock deviation in aged
individuals.
\end{corollary}

%------------------------------------------------------------------------------
\section{Phase Separation as Categorical Partition Phase Transition}
\label{sec:phase_separation}
%------------------------------------------------------------------------------

Liquid-liquid phase separation (LLPS) in the nucleus has received enormous
attention as an organising principle for nuclear bodies (nucleolus, PML
bodies, Cajal bodies), heterochromatin domains, and transcription hubs. The
biophysical description treats phase separation as a standard polymer
phase transition driven by multivalent weak interactions above a
critical concentration.

Partition theory reinterprets phase separation as a \emph{categorical partition
phase transition} --- a switch between two distinct categorical states, not a
continuum of intermediate states.

\begin{theorem}[Phase Separation as Categorical Phase Transition]
\label{thm:phase_separation}
Chromatin phase separation occurs at a critical partition depth $\Depth_{\mathrm{crit}}$:
\begin{enumerate}[label=(\roman*)]
  \item For $\Depth(x) < \Depth_{\mathrm{crit}}$: chromatin at locus $x$ is
        in the \emph{euchromatin phase} (dilute, accessible, active).
  \item For $\Depth(x) > \Depth_{\mathrm{crit}}$: chromatin at locus $x$ is
        in the \emph{heterochromatin phase} (condensed, inaccessible, silent).
\end{enumerate}
The critical depth $\Depth_{\mathrm{crit}}$ is determined by the balance
\begin{equation}
  \Depth_{\mathrm{crit}} = \frac{\ln([\mathrm{HP1}]/K_{d,\mathrm{HP1}})}{\ln b}
  + \Delta\Depth_{\mathrm{H3K9me3}},
  \label{eq:phase_crit_depth}
\end{equation}
where $[\mathrm{HP1}]$ is the local HP1 concentration, $K_{d,\mathrm{HP1}}$
is the HP1 chromodomain binding affinity for H3K9me3, and
$\Delta\Depth_{\mathrm{H3K9me3}}$ is the depth increment of H3K9me3
(from Theorem~\ref{thm:histone_code} in Chapter~\ref{chap:chromatin_splicing}).
\end{theorem}

\begin{proof}
HP1 (CBX5 in mammals) binds H3K9me3 through its chromodomain
($K_d \approx 1$--$10\,\mu$M) and dimerises through its chromoshadow domain.
At concentrations above $K_{d,\mathrm{HP1}}$, HP1 bridges adjacent
H3K9me3-marked nucleosomes, increasing the effective partition depth of the
region through steric occlusion and recruitment of additional silencing factors.
The partition depth evolution equation is:
$\dot\Depth = k_{\mathrm{HP1}}[\mathrm{HP1}]/\Depth - k_{\mathrm{loss}}\Depth$.
The fixed point $\Depth^* = \sqrt{k_{\mathrm{HP1}}[\mathrm{HP1}]/k_{\mathrm{loss}}}$
bifurcates at $[\mathrm{HP1}] = K_{d,\mathrm{HP1}} \exp(\Depth_{\mathrm{crit}}\ln b)$,
which rearranges to Eq.~\eqref{eq:phase_crit_depth}. Below $\Depth_{\mathrm{crit}}$,
the system has a single low-$\Depth$ fixed point (euchromatin); above it,
the system has a high-$\Depth$ fixed point (heterochromatin). This is a
categorical (discontinuous) bifurcation, not a continuous crossover. $\blacksquare$
\end{proof}

\begin{remark}
Phase separation in the biophysical sense (polymer demixing, driven by
multivalent interactions) provides the physical mechanism by which the
categorical partition transition is realised in three dimensions. The partition
theory provides the \emph{criterion} (depth threshold $\Depth_{\mathrm{crit}}$),
while biophysics provides the \emph{mechanism} (HP1-driven polymer demixing).
The two descriptions are complementary, operating at different levels of
abstraction.
\end{remark}

%------------------------------------------------------------------------------
\section{Heterochromatin as Charge-Isolated Partition}
\label{sec:heterochromatin_partition}
%------------------------------------------------------------------------------

\begin{proposition}[Heterochromatin as Maximal-Depth Partition]
\label{prop:heterochromatin_max_depth}
Constitutive heterochromatin (pericentromeric and telomeric regions, marked
by H3K9me3 + HP1) represents chromatin at or near $\Depth_{\mathrm{max}}$
for the genomic partition. At $\Depth = \Depth_{\mathrm{max}}$:
\begin{enumerate}[label=(\roman*)]
  \item Transcription factor access approaches zero (partition gate fully
        closed, $\theta \approx 0$).
  \item Replication firing time $T_{\mathrm{rep}}$ is maximal (last to
        replicate in S-phase, consistent with heterochromatin being
        constitutively late-replicating; Theorem~\ref{thm:replication_timing}
        in Chapter~\ref{chap:metabolic_clock}).
  \item Electrostatic charge is maximally screened by HP1 dimers, condensins,
        and \ce{Mg^{2+}}: the Debye length oscillation has minimal amplitude
        effect at $\Depth_{\mathrm{max}}$ because the region is already fully
        collapsed.
\end{enumerate}
Heterochromatin is the genome's charge-isolated partition: a region sequestered
from the oscillating nuclear electrostatic field by its own depth.
\end{proposition}

The H3K9me3/HP1 pathway is therefore a charge-isolation mechanism: it builds
partition depth not to regulate gene expression at that locus (constitutive
heterochromatin contains very few protein-coding genes) but to physically
isolate the underlying sequence (repetitive elements, satellite DNA) from the
transcriptional and replicative machinery, preventing genomic instability
from transposable element activation and repeat-mediated recombination.

%------------------------------------------------------------------------------
\section{The Centromere as the Primary Partition Origin of the Chromosome}
\label{sec:centromere_partition}
%------------------------------------------------------------------------------

Each chromosome has a centromere --- a specialised locus at which the
kinetochore is assembled during mitosis and to which spindle microtubules
attach. The centromere is defined not by a specific DNA sequence (in most
organisms) but by the presence of CENP-A-containing nucleosomes, in which
the canonical H3 is replaced by the centromere-specific histone variant CENP-A.

\begin{definition}[Centromere as Partition Origin]
\label{def:centromere_partition_origin}
The centromere of a chromosome is the partition origin: the unique locus with
the deepest partition depth on the chromosome, from which the chromosome's
physical categorical identity --- as a distinct unit to be segregated ---
is computed. Its depth exceeds that of any other locus on the chromosome:
\begin{equation}
  \Depth_{\mathrm{centromere}} > \Depth(x) \quad \forall\, x \neq \mathrm{centromere}.
  \label{eq:centromere_maxdepth}
\end{equation}
\end{definition}

\begin{theorem}[CENP-A Nucleosome Marks the Deepest Partition]
\label{thm:cenpa_deepest}
CENP-A nucleosomes have a greater partition depth than canonical H3 nucleosomes:
\begin{enumerate}[label=(\roman*)]
  \item CENP-A has a more positively charged histone fold domain than H3,
        increasing $|U_{\mathrm{bind}}|$ (Eq.~\eqref{eq:nucleosome_binding_energy})
        and reducing the breathing angle $\theta$ to near-zero.
  \item CENP-A nucleosomes recruit CCAN (Constitutive Centromere-Associated
        Network) proteins, each of which adds further positive charges that
        deepen the electrostatic well, increasing $\Depth_{\mathrm{centromere}}$.
  \item The CENP-A chromatin domain is therefore the lowest-breathing, highest-depth
        chromatin on the chromosome: a stable categorical landmark that does
        not oscillate with the metabolic clock, providing a constant reference
        point for kinetochore assembly.
\end{enumerate}
\end{theorem}

\begin{proof}
Compare the breathing equation (Theorem~\ref{thm:nucleosome_breathing}) for
CENP-A and H3 nucleosomes. CENP-A has 6 additional basic residues in its
C-terminal tail compared to H3, increasing $Q_{\mathrm{histones}}$ by
$\sim 6e$ per CENP-A copy, or $\sim 12e$ per nucleosome. From
Eq.~\eqref{eq:nucleosome_binding_energy}, this increases $|U_{\mathrm{bind}}|$
by $\Delta U = 12e \cdot Q_{\mathrm{DNA}} / (4\pi\varepsilon_0\varepsilon_r r_{\mathrm{eff}})
\approx 2.5\,\kB T$ per nucleosome (using $r_{\mathrm{eff}} = 2$\,nm).
From Theorem~\ref{thm:nucleosome_breathing}, the breathing amplitude decreases
by $\delta\theta/\bar\theta = -\Delta U/\kB T \approx -2.5$, driving
$\theta \to 0$. With $\theta \approx 0$, $\Depth = \log_b(\Omega_{\rm total}/(\theta\,\Omega_{\rm open})) \to \infty$:
the centromere nucleosome is a permanently closed partition gate. $\blacksquare$
\end{proof}

%------------------------------------------------------------------------------
\section{Nuclear Bodies as Partition Condensates}
\label{sec:nuclear_bodies}
%------------------------------------------------------------------------------

The nucleus contains multiple non-membrane-bound compartments: the nucleolus,
Cajal bodies, PML bodies, splicing speckles, and transcription hubs (also
called super-enhancer condensates). Each is a phase-separated condensate.
Partition theory classifies them as partition condensates formed at specific
depth thresholds.

\begin{definition}[Partition Condensate Hierarchy]
\label{def:condensate_hierarchy}
The nuclear partition condensates are ordered by their characteristic depth
threshold $\Depth_{\mathrm{crit}}^{(i)}$:
\begin{center}
\begin{tabular}{lll}
\toprule
Condensate & Primary driver & $\Depth_{\mathrm{crit}}$ (relative) \\
\midrule
Transcription hubs & Med1, BRD4 / active enhancers & lowest \\
Splicing speckles & SRSF1/2, SC35 / active pre-mRNA & low \\
Cajal bodies & Coilin, SMN / snRNA biogenesis & medium \\
Nucleolus & NPM1, FBL / rDNA transcription & medium \\
PML bodies & PML, SP100 / DNA repair & high \\
Constitutive heterochromatin & HP1, H3K9me3 / repeats & highest \\
\bottomrule
\end{tabular}
\end{center}
Each condensate is a region where chromatin or RNA reaches its characteristic
$\Depth_{\mathrm{crit}}$ and undergoes the categorical partition phase
transition of Theorem~\ref{thm:phase_separation}.
\end{definition}

The ordering in Definition~\ref{def:condensate_hierarchy} is not arbitrary:
it reflects the functional hierarchy of the genome. Transcription hubs have
the lowest $\Depth_{\mathrm{crit}}$ because they must be most responsive to
metabolic oscillations (they form and dissolve rapidly with the transcriptional
burst cycle). Constitutive heterochromatin has the highest $\Depth_{\mathrm{crit}}$
because it must be maximally stable and resistant to metabolic fluctuations.

%------------------------------------------------------------------------------
\section{Unified Picture: Nuclear Architecture as Partition Hierarchy}
\label{sec:unified_nuclear}
%------------------------------------------------------------------------------

The preceding sections establish that every structural level of nuclear
organisation corresponds to a level in the chromatin partition hierarchy.

\begin{proposition}[Nuclear Architecture as Realised Partition Hierarchy]
\label{prop:nuclear_architecture_hierarchy}
The nuclear architecture of a eukaryotic cell implements the following
partition hierarchy, from coarsest to finest:
\begin{align}
  &\Part_{\mathrm{cell}} && (\text{nucleus vs.\ cytoplasm, NPC gate})
  \nonumber\\
  &\quad\hookrightarrow \Part_{\mathrm{chromosome}}^{(k)} && (\text{chromosome territories, charge-minimised packing})
  \nonumber\\
  &\quad\quad\hookrightarrow \Part_{\mathrm{TAD}}^{(k,l)} && (\text{charge domains, TAD structure})
  \nonumber\\
  &\quad\quad\quad\hookrightarrow \Part_{\mathrm{loop}}^{(k,l,m)} && (\text{chromatin loops, cohesin-extruded})
  \nonumber\\
  &\quad\quad\quad\quad\hookrightarrow \Part_{\mathrm{nucleosome}}^{(k,l,m,n)} && (\text{nucleosome gates, breathing control})
  \nonumber\\
  &\quad\quad\quad\quad\quad\hookrightarrow \Part_{\mathrm{bp}}^{(k,l,m,n,p)} && (\text{base pair sequence, TF binding sites})
\end{align}
Each level $\hookrightarrow$ is a partition refinement: each higher level
partitions the elements of the level above into finer sub-partitions.
The metabolic clock drives oscillations simultaneously at all levels,
with amplitude decreasing from the nucleosome level upward (higher-level
partitions have larger effective $\Depth$ and are more damped) and from
the centromere outward (the centromere anchors the deepest partition).
\end{proposition}

%------------------------------------------------------------------------------
\section{Quantitative Predictions}
\label{sec:architecture_predictions}
%------------------------------------------------------------------------------

\begin{proposition}[NPC Depth Degradation is Detectable by FG-Nup FRAP]
\label{prop:npc_frap}
If NPC partition depth decreases with age according to Eq.~\eqref{eq:npc_damage},
then the fluorescence recovery after photobleaching (FRAP) half-time of
inert fluorescent dextrans above the 40\,kDa cutoff should decrease with
donor age in neurons. Conversely, the FRAP half-time of 10\,kDa dextrans
(freely passing) should be independent of age.
\end{proposition}

\begin{proposition}[CENP-A Breathing Amplitude is Below Detection Threshold]
\label{prop:cenpa_breathing}
From Theorem~\ref{thm:cenpa_deepest}, the breathing angle of CENP-A
nucleosomes satisfies $\theta < 0.01$ under physiological conditions.
Single-molecule FRET measurements should be unable to detect breathing
oscillations at centromeric loci at metabolic timescales, while pericentromeric
H3 nucleosomes (with $\theta \approx 0.3$--$0.5$) should show clear metabolic-
frequency oscillations. The boundary between breathing and non-breathing
nucleosomes should coincide with the CENP-A/H3 boundary detected by CENP-A
ChIP-seq.
\end{proposition}

\begin{proposition}[Phase Separation Threshold Scales with HP1 Concentration]
\label{prop:phase_sep_hp1}
From Eq.~\eqref{eq:phase_crit_depth}, the volume fraction of heterochromatin
in a cell should scale as a sigmoidal function of HP1 (CBX5) concentration
with inflection point at $[\mathrm{HP1}] = K_{d,\mathrm{HP1}}$. Overexpression
of HP1 beyond this point should drive additional euchromatin into
heterochromatin, reducing $\Se$ (evolutionary entropy) but also reducing
global transcription rate (increased $\Depth_{\mathrm{crit}}$ suppresses more
promoters). This prediction is testable by ChIP-seq and RNA-seq in HP1-
overexpressing cells.
\end{proposition}

%------------------------------------------------------------------------------
\begin{chapterbox}
\textbf{Chapter Summary.}
\begin{itemize}
  \item The nuclear envelope defines the highest-level categorical partition
        of the eukaryotic cell. The NPC implements a partition filter with
        depth $\Depth_{\mathrm{filter}}(m) = \alpha m$, passing small
        molecules freely and excluding large macromolecules
        (Theorem~\ref{thm:npc_filter}).
  \item Aging is partition boundary dissolution: NPC damage accumulates
        exponentially (Theorem~\ref{thm:aging_npc}), progressively
        reducing $\Depth_{\mathrm{filter}}$ and causing inappropriate
        nuclear--cytoplasmic exchange. This is the molecular mechanism of
        age-dependent epigenetic drift (Corollary~\ref{cor:epigenetic_drift}).
  \item Chromatin phase separation is a categorical partition phase
        transition at depth $\Depth_{\mathrm{crit}}$
        (Theorem~\ref{thm:phase_separation}): euchromatin
        ($\Depth < \Depth_{\mathrm{crit}}$) and heterochromatin
        ($\Depth > \Depth_{\mathrm{crit}}$) are two discrete categorical states,
        not endpoints of a continuum.
  \item The centromere is the deepest partition of the chromosome
        (Definition~\ref{def:centromere_partition_origin}): CENP-A
        nucleosomes have $\theta \approx 0$ (fully closed partition gate),
        providing the stable categorical landmark for kinetochore assembly
        (Theorem~\ref{thm:cenpa_deepest}).
  \item Nuclear architecture is a realised partition hierarchy from cell
        boundary to base pair, with each level a partition refinement of
        the level above (Proposition~\ref{prop:nuclear_architecture_hierarchy}).
  \item The oscillating nuclear electric field ($\sim 10^5$\,V\,m$^{-1}$,
        established in Chapter~\ref{chap:transcriptional_bursting}) drives
        all partition levels simultaneously, with the metabolic clock
        as the master oscillator.
\end{itemize}
\end{chapterbox}
